%% kapa: exam=1 round=30 year=2019 period=1 form=A subject=civil_law subject_ko=민법 seq=3 number=3
%% source: Q-Net 감정평가사 기출문제 (cst00309) artlSeq=5205205 file=extracted/1008/5205205/2019년 제30회 감정평가사 1차시험 문제지/2019년 제30회 감정평가사 1차시험 문제지/2019년도 제30회 감정평가사 1차 1교시 A형.hwp
%% conversion: kordoc 4.24.51 -> tools/gen1.py (auto); review=proofread-sonnet; flags=none
\begin{question}{3}
신의칙과 권리남용에 관한 설명으로 옳지 않은 것은? (다툼이 있으면 판례에 따름)
\begin{choices}
  \item 신의칙에 반하는 것인지 여부는 당사자의 주장이 없더라도 법원이 직권으로 판단할 수 있다.
  \item 신의칙에 기한 사정변경의 원칙에 의하여 계약해제권이 발생할 수 있다.
  \item 강행법규에 반한다는 사정을 알면서 법률행위를 한 자가 강행법규 위반을 이유로 그 법률행위의 무효를 주장하는 것은 특별한 사정이 없는 한 신의칙에 위배되지 않는다.
  \item 권리남용금지의 원칙은 본래적 의미의 권리뿐만 아니라 법인격의 남용에도 적용된다.
  \item 국민을 보호할 의무가 있는 국가가 국민에 대하여 부담하는 손해배상채무의 소멸시효 완성을 주장하는 것은 원칙적으로 신의칙에 반한다.
\end{choices}
\end{question}
\answer{5}
